\begin{table}[H]
\centering\small
\begin{tabular}{lcccc}
\toprule
Estrategia & ROC-AUC & PR-AUC & Brier & pend. cal. \\
\midrule
\textbf{E3 deslizante S=3} & 0.680 & 0.226 & 0.095 & 1.123 \\
E2 expansiva & 0.676 & 0.223 & 0.095 & 1.075 \\
E3 deslizante S=2 & 0.676 & 0.222 & 0.095 & 1.053 \\
E4 disparador & 0.674 & 0.222 & 0.095 & 1.166 \\
E3 deslizante S=5 & 0.679 & 0.219 & 0.095 & 1.039 \\
E3 deslizante S=1 & 0.667 & 0.214 & 0.095 & 0.918 \\
E1 estático & 0.662 & 0.212 & 0.097 & 1.028 \\
\bottomrule
\end{tabular}
\caption{Comparación de estrategias de adaptación sobre bloques comunes (6--10). \textbf{E3 deslizante $S=3$} obtiene el mejor PR-AUC promedio entre E1/E2/E3 y se elige como estrategia desplegada para la capa de decisión --- coincide con la ventana de reentrenamiento de tamaño 3 propuesta originalmente en el plan de trabajo. E4 (disparador) se muestra por separado por ser condicional a los reentrenamientos que dispara (1 evento en el horizonte evaluado).}
\label{tab:estrategias}
\end{table}
